\chapter{源码阅读、Benchmark Suite 与工程报告}

前面章节讲的是原理和局部内核。真正进入工程时，读者会面对复杂代码库：平台宏、模板、代码生成、线程池、内存池、图优化、量化格式和大量 benchmark。源码阅读不能靠从第一行读到最后一行，而要带着问题建立地图。本章给出一套阅读和验证方法。

\section{先找数据结构}

高性能 AI 项目的核心通常不是某个孤立函数，而是一组数据结构。矩阵或张量如何描述 shape、stride、dtype、layout；权重如何打包；线程任务如何表示；KV Cache 如何分页；量化 block 如何保存 scale；benchmark 如何描述输入形状。这些结构决定后续函数能做什么。

阅读 oneDNN 时，可以先找 primitive descriptor、memory descriptor 和具体 ISA kernel 的关系。阅读 OpenBLAS 或 BLIS 时，可以先找 GEMM 的 block size、packing buffer 和 microkernel 表。阅读 XNNPACK 时，可以先找 operator、microkernel、indirection buffer 和 threadpool。阅读 llama.cpp 时，可以先找 tensor、ggml graph、backend、KV cache 和量化 block。

\section{再找热路径}

热路径是实际运行中占主要时间的代码。可以通过 profiler 找，也可以从模型结构推断。大语言模型推理中，矩阵乘、Attention、RMSNorm、RoPE、采样和 KV Cache 操作都可能在热路径上。读源码时要区分热路径和配置路径。配置路径复杂但只执行一次，热路径哪怕一行也可能影响吞吐。

阅读热路径时建议画出数据流：输入张量从哪里来，经过哪个 layout，是否 packing，进入哪个 kernel，输出写到哪里，是否立即被下一个算子读取。每个边都标注 dtype、shape 和内存位置。这样做比单纯抄函数调用栈更有用，因为性能瓶颈通常发生在边上，也就是数据移动处。

\section{Benchmark Suite 的结构}

单个 benchmark 容易片面。一个工程化 benchmark suite 应该覆盖多个层级：微基准测单个循环或内核，算子基准测 GEMM、卷积、Softmax、Attention，模型片段基准测一层 Transformer，端到端基准测 prefill、decode、tokens/s 和 latency。每个层级回答不同问题。

基准测试还要覆盖形状。只测一个大矩阵会让 GEMM 看起来很好，却无法解释小 batch decode；只测短序列无法解释长上下文 KV Cache；只测平均值无法解释服务尾延迟。一个合理 suite 至少包含大中小矩阵、不同 batch、不同 sequence length、不同线程数、不同 dtype 和冷/热缓存场景。

\begin{lstlisting}[language=C++,caption={Benchmark case 描述应显式记录形状和线程数}]
#include <cstdint>
#include <string>

struct BenchmarkCase {
    std::string name;
    std::int32_t batch;
    std::int32_t sequence_length;
    std::int32_t hidden_size;
    std::int32_t heads;
    std::int32_t head_dim;
    std::int32_t threads;
};
\end{lstlisting}

记录结构化 case 的好处是，结果可以长期比较。后续如果改了 packing、线程切分或量化格式，可以直接看哪些形状变快，哪些形状变慢。没有结构化 case，benchmark 只是一堆难以追溯的数字。

\section{工程报告怎么写}

性能报告应当先写结论，再写证据。结论包括优化目标、变更范围、主要收益和风险。证据包括硬件环境、编译选项、输入形状、测量方法、统计方式、性能表、计数器和正确性验证。报告还要写失败尝试，因为失败尝试能说明边界，避免后续重复踩坑。

一份合格报告不应该只说“GEMM 优化后快了 30\%”。它要说明在哪个 CPU、什么矩阵形状、什么线程数、什么 dtype、基线是谁、重复几次、误差如何、是否影响其他形状、为什么快。若推测原因是 cache miss 降低，就给出 perf 证据；若推测原因是向量化成功，就给出汇编或编译器报告。

\begin{keyidea}
源码阅读、benchmark 和报告是一件事的三个面：源码给出机制，benchmark 给出事实，报告把机制和事实连接成可复查的工程结论。
\end{keyidea}

\section{后续扩写方向}

本册后续可以继续扩写四类材料。第一类是更完整的实验代码，包括 GEMM 分块、AVX2 点积、Softmax 优化、int8 量化和线程池。第二类是源码专题，逐步拆解 OpenBLAS、oneDNN、XNNPACK、llama.cpp 和 Linux 调度/内存相关代码。第三类是模型案例，例如 MLP、Attention、KV Cache 长上下文和端到端 decode。第四类是报告模板和性能回归工具，让每次优化都能留下证据。

\begin{exercisebox}
选择一个开源项目中的 GEMM 或 Attention benchmark，整理它的输入形状、线程设置、计时范围和正确性检查。指出它能回答什么问题，不能回答什么问题，并提出三个应该补充的 case。
\end{exercisebox}
